\section{Analysis: why averaging can help}
\label{sec:analysis}
Averaging can improve predictions when spatial errors partly offset. For a
fixed library, classical ensemble analysis writes the expected squared relative
error as
\begin{equation}
 R(w)=\E[w^\top G(x,y)w]=w^\top Cw,\qquad C=\E[G(x,y)].
 \label{eq:riskmain}
\end{equation}
Here $G(x,y)$ is the error product matrix for a new input and target pair,
$\E$ averages over such pairs, $C$ is its population mean, and $R(w)$ is the
mixture's expected squared relative error. Its diagonal measures individual
error and its other entries retain alignment and shared bias
\citep{breiman1996,krogh1994}. The inverse error mixture can exploit complementary
predictions even though it uses only individual accuracy to choose weights.

This identity does not guarantee useful weights from a small fitting set,
and the bound in Appendix~\ref{app:proof} does not certify accuracy from five fields.
The reported metric $L(w)=\E[\sqrt{w^\top Gw}]$ averages unsquared relative error
and can favor different weights from $R$ even with unlimited fitting data.
Appendix~\ref{app:geometry} gives a counterexample and a direct fitting control.
